\documentclass[aspectratio=169]{beamer}

\usepackage[utf8]{inputenc}
\usepackage[T1]{fontenc}
\usepackage[french]{babel}
\usepackage{amsmath, amssymb}
\usepackage{booktabs}

\usetheme{Madrid}
\setbeamertemplate{navigation symbols}{}

\title[Modèle HistGradientBoosting]{Prévision de la consommation des accouplements\\Modèle HistGradientBoostingRegressor}
\author{Projet OCP -- Système de prévision}
\date{}

\begin{document}

\begin{frame}
  \titlepage
\end{frame}

%--------------------------------------------------------------------
\begin{frame}{Objectif}
  \begin{itemize}
    \item Prévoir la consommation annuelle 2025 de chaque accouplement.
    \item En déduire une prévision mensuelle $n_{pm}$ et une quantité à commander $n_{ac}$ pour le mois suivant (M+1).
    \item Modèle retenu : \textbf{HistGradientBoostingRegressor} (\textsc{scikit-learn}).
  \end{itemize}
\end{frame}

%--------------------------------------------------------------------
\begin{frame}{Jeu de données}
  \begin{itemize}
    \item Panel de consommation annuelle par pièce (2012--2024).
    \item Référentiel pièces (\texttt{part2\_clean\_step2.parquet}) :
      \begin{itemize}
        \item informations de stock (stock actuel, stock de sécurité, min/max),
        \item quantité installée, demande cumulée 3 ans.
      \end{itemize}
    \item Une ligne d'apprentissage = \textbf{une pièce} avec ses variables dérivées.
  \end{itemize}
\end{frame}

%--------------------------------------------------------------------
\begin{frame}{Features (1/3) -- Lags annuels}
  \textbf{Historique de consommation par pièce :}
  \begin{itemize}
    \item \textbf{lag1} : consommation annuelle année $N-1$.
    \item \textbf{lag2} : consommation annuelle année $N-2$.
    \item \textbf{lag3} : consommation annuelle année $N-3$.
  \end{itemize}
  \vspace{0.5em}
  \begin{block}{Rôle}
    Ces lags capturent le niveau récent et les variations à court terme de la consommation.
  \end{block}
\end{frame}

%--------------------------------------------------------------------
\begin{frame}{Features (2/3) -- Statistiques sur l'historique}
  \begin{itemize}
    \item \textbf{cons\_mean\_3y} : moyenne des 3 dernières années
      \[
        \text{cons\_mean\_3y} = (\text{lag1} + \text{lag2} + \text{lag3})/3.
      \]
    \item \textbf{cons\_trend} : tendance récente
      \[
        \text{cons\_trend} = \text{lag1} - \text{lag2}.
      \]
    \item \textbf{cons\_mean\_all\_past} : moyenne de toutes les années historiques.
  \end{itemize}
  \vspace{0.5em}
  \begin{block}{Rôle}
    Résumer le niveau moyen et la pente de long terme de la consommation.
  \end{block}
\end{frame}

%--------------------------------------------------------------------
\begin{frame}{Features (3/3) -- Structure et valeurs manquantes}
  \textbf{Caractéristiques structurelles de la pièce :}
  \begin{itemize}
    \item \textbf{qte\_installee} (parc installé).
    \item \textbf{stock\_actuel}, \textbf{stock\_securite}, \textbf{stock\_min}, \textbf{stock\_max}.
    \item \textbf{qte\_demandee\_3ans} (demande totale sur 3 ans).
  \end{itemize}
  \vspace{0.5em}
  \textbf{Indicateurs de données manquantes :}
  \begin{itemize}
    \item \textbf{lag1\_was\_missing}, \textbf{lag2\_was\_missing}, \textbf{lag3\_was\_missing}.
    \item \textbf{qte\_installee\_was\_missing}, \textbf{stock\_securite\_was\_missing}.
  \end{itemize}
  \begin{block}{Rôle}
    Intégrer la criticité opérationnelle de la pièce et informer le modèle lorsqu'une donnée est imputée.
  \end{block}
\end{frame}

%--------------------------------------------------------------------
\begin{frame}{Pourquoi HistGradientBoosting ?}
  \begin{itemize}
    \item \textbf{Non-linéarités et interactions} :
      \begin{itemize}
        \item relation complexe entre lags, stocks et demandes,
        \item les arbres de décision + boosting capturent naturellement ces effets.
      \end{itemize}
    \item \textbf{Gestion des distributions skewed} :
      \begin{itemize}
        \item beaucoup de pièces à faible consommation, quelques pièces très consommatrices,
        \item les arbres sont robustes aux valeurs extrêmes.
      \end{itemize}
    \item \textbf{Implémentation efficace} :
      \begin{itemize}
        \item version histogramme de scikit-learn, bien adaptée aux features numériques,
        \item bon compromis vitesse / performance.
      \end{itemize}
    \item \textbf{Par rapport au Random Forest} :
      \begin{itemize}
        \item RF obtient de bonnes performances mais a tendance à sur-apprendre sur certaines pièces rares,
        \item le boosting, avec une profondeur limitée et un taux d'apprentissage modéré, est plus régularisé et
              donne un meilleur MAE/RMSE sur le jeu de validation.
      \end{itemize}
    \item \textbf{Meilleures performances sur validation} que les baselines, Ridge et Random Forest.
  \end{itemize}
\end{frame}

%--------------------------------------------------------------------
\begin{frame}{Procédure d'entraînement}
  \begin{itemize}
    \item Sélection des années complètes 2012--2024.
    \item Construction du jeu de features par pièce.
    \item Découpage des données :
      \begin{itemize}
        \item \textbf{train} : apprentissage des paramètres,
        \item \textbf{validation} : choix du modèle / hyperparamètres,
        \item \textbf{test} : évaluation finale.
      \end{itemize}
    \item Hyperparamètres du modèle retenu :
      \begin{itemize}
        \item \texttt{learning\_rate} = 0.1, \texttt{max\_depth} = 3,
        \item \texttt{max\_iter} = 200, \texttt{min\_samples\_leaf} = 5.
      \end{itemize}
  \end{itemize}
\end{frame}

%--------------------------------------------------------------------
\begin{frame}{Métriques d'évaluation}
  \begin{itemize}
    \item \textbf{MAE} : Mean Absolute Error.
    \item \textbf{RMSE} : Root Mean Squared Error.
    \item \textbf{ACC\_within\_1} : proportion de prévisions dont l'erreur absolue $\leq 1$ unité.
  \end{itemize}
  \vspace{0.5em}
  Modèles comparés :
  \begin{itemize}
    \item baselines (\textit{last value}, \textit{mean 3y}),
    \item Ridge,
    \item Random Forest,
    \item HistGradientBoosting.
  \end{itemize}
\end{frame}

%--------------------------------------------------------------------
\begin{frame}{Résultats comparés (validation / test)}
  \small
  \begin{tabular}{lcccc}
    \toprule
    Modèle & Split & MAE & RMSE & ACC\_within\_1 \\
    \midrule
    Baseline last value & val & 0.80 & 1.44 & 0.79 \\
                        & test & 0.68 & 1.74 & 0.83 \\
    Baseline mean 3y    & val & 0.73 & 1.39 & 0.81 \\
                        & test & 0.73 & 1.70 & 0.82 \\
    Ridge               & val & 0.69 & 1.20 & 0.80 \\
                        & test & 0.69 & 1.44 & 0.80 \\
    Random Forest       & val & 0.44 & 0.80 & 0.88 \\
                        & test & 0.72 & 1.46 & 0.78 \\
    \textbf{HistGradBoost} & val & \textbf{0.40} & \textbf{0.72} & \textbf{0.90} \\
                        & test & 0.76 & 1.55 & 0.78 \\
    \bottomrule
  \end{tabular}
\end{frame}

%--------------------------------------------------------------------
\begin{frame}{Conclusion}
  \begin{itemize}
    \item Ensemble compact de features :
      \begin{itemize}
        \item lags, moyennes, tendance,
        \item caractéristiques de stock et indicateurs de données manquantes.
      \end{itemize}
    \item HistGradientBoosting offre le meilleur compromis :
      \begin{itemize}
        \item erreurs (MAE / RMSE) les plus faibles sur validation,
        \item modèle robuste et interprétable pour un usage opérationnel.
      \end{itemize}
    \item Les prévisions obtenues alimentent le calcul de $n_{pm}$ et de la quantité à commander $n_{ac}$ pour la gestion des accouplements.
  \end{itemize}
\end{frame}

\end{document}
